% year: תשפ"ג
% semester: ב
% moed: א
% source: exams/מבחנים/תשפג/מועד א תשפג סמסטר ב.pdf
% number: 4א
% points: 20
% categories: lhopital, func-limits

\begin{question}
חשבו את הגבולות הבאים
\[ (1)\ \lim_{x\to0}\frac{\ln(1+2x)}{\sin2x-\sin6x}\qquad\qquad (2)\ \lim_{x\to0}\left(\sqrt{1+2x}\right)^{\frac{1}{(-x)}} \]
\end{question}

\begin{hint}
בגבול (2) רשמו את הביטוי בצורה $e^{g(x)}$ וחשבו את הגבול של המעריך.
\end{hint}

\begin{solution}
\textbf{(1)} גבול מהצורה $\frac00$; לפי כלל לופיטל:
\[ \lim_{x\to0}\frac{\ln(1+2x)}{\sin2x-\sin6x}=\lim_{x\to0}\frac{\frac{2}{1+2x}}{2\cos2x-6\cos6x}=\frac{2}{2-6}=\boxed{-\frac12}. \]
(לחלופין, בעזרת $\ln(1+2x)\sim2x$, $\sin2x\sim2x$, $\sin6x\sim6x$: $\frac{2x}{2x-6x}=-\frac12$.)

\textbf{(2)} עבור $x$ קרוב ל-$0$ ($x>-\frac12$, $x\neq0$) הבסיס חיובי, ו-
\[ \left(\sqrt{1+2x}\right)^{-1/x}=e^{-\frac{\ln(1+2x)}{2x}}. \]
לפי הגבול היסודי (או לופיטל) $\lim_{x\to0}\frac{\ln(1+2x)}{2x}=1$. מרציפות פונקציית האקספוננט,
\[ \lim_{x\to0}\left(\sqrt{1+2x}\right)^{\frac{1}{-x}}=e^{-1}=\boxed{\frac1e}. \]
\end{solution}
